\documentclass[10pt,letterpaper,oneside]{report}
\usepackage[T1]{fontenc}
\usepackage[utf8]{inputenc}
\usepackage{lmodern}
\usepackage{microtype}
\usepackage[margin=0.85in,headheight=15pt]{geometry}
\usepackage{enumitem,booktabs,longtable,tabularx,array,makecell}
\usepackage{fancyhdr,lastpage,xcolor,hyperref,titlesec}
\usepackage{listings}
\usepackage[most]{tcolorbox}
\definecolor{ink}{HTML}{182D40}
\definecolor{accent}{HTML}{126B77}
\definecolor{pale}{HTML}{EEF5F5}
\definecolor{grayblue}{HTML}{677887}
\definecolor{codebg}{HTML}{F5F7F8}
\hypersetup{colorlinks=true,linkcolor=accent,urlcolor=accent,pdfauthor={C23 Engineering Companion},pdfsubject={ISO/IEC 9899:2024 educational engineering companion}}
\titleformat{\chapter}[display]{\normalfont\huge\bfseries\color{ink}}{\filleft\large\color{accent}CHAPTER \thechapter}{1ex}{\titlerule[1pt]\vspace{1ex}}
\titleformat{\section}{\Large\bfseries\color{ink}}{\thesection}{0.7em}{}
\titleformat{\subsection}{\large\bfseries\color{accent}}{\thesubsection}{0.65em}{}
\setlist[itemize]{leftmargin=1.3em,itemsep=2pt,topsep=3pt}
\setlist[enumerate]{leftmargin=1.6em,itemsep=3pt,topsep=3pt}
\setlength{\parskip}{5pt}
\setlength{\parindent}{0pt}
\renewcommand{\arraystretch}{1.18}
\pagestyle{fancy}\fancyhf{}\fancyhead[L]{\small\sffamily C23 ENGINEERING SERIES}\fancyhead[R]{\small\sffamily\color{accent}\DocShort}\fancyfoot[L]{\small\color{grayblue}ISO/IEC 9899:2024 companion}\fancyfoot[R]{\small\thepage\ / \pageref{LastPage}}
\newtcolorbox{keybox}{breakable,colback=pale,colframe=accent,boxrule=.6pt,arc=1mm,left=8pt,right=8pt,top=6pt,bottom=6pt}
\newtcolorbox{warnbox}{breakable,colback=codebg,colframe=grayblue,boxrule=.6pt,arc=1mm,left=8pt,right=8pt,top=6pt,bottom=6pt}
\lstdefinestyle{cstyle}{language=C,basicstyle=\ttfamily\footnotesize,keywordstyle=\color{accent}\bfseries,commentstyle=\color{grayblue}\itshape,stringstyle=\color{ink},backgroundcolor=\color{codebg},frame=single,rulecolor=\color{grayblue!40},breaklines=true,showstringspaces=false,tabsize=4,columns=fullflexible,keepspaces=true,upquote=true}
\lstset{style=cstyle}
\newcommand{\refcl}[1]{\textsf{\footnotesize[ISO/IEC 9899:2024, #1]}}
\newcommand{\standard}{\textbf{[STANDARD]}}
\newcommand{\optional}{\textbf{[OPTIONAL]}}
\newcommand{\guidance}{\textbf{[GUIDANCE]}}
\newcommand{\testtag}{\textbf{[TEST]}}
\newcommand{\term}[1]{\textbf{#1}}
\newcommand{\DocShort}{CURRICULUM}
\begin{document}
\begin{titlepage}
\color{ink}\vspace*{1.2in}
{\sffamily\small\color{accent}\bfseries TECHNICAL DOCUMENTATION SUITE \hfill EDITION 2026-10}\par
\vspace{.3in}\rule{\linewidth}{1.5pt}\par\vspace{.22in}
{\Huge\bfseries C23 Prerequisite-Based Training Curriculum}\par\vspace{.22in}
{\Large\color{grayblue}18-week technical progression, assessments and capstone gates}\par\vspace{.4in}
\begin{keybox}\textbf{Standards basis:} ISO/IEC 9899:2024 (C23). Independent educational and engineering companion. Not an official ISO interpretation, full specification, certification or conformance assessment.\end{keybox}
\vfill{\large\bfseries C23 Software Engineering and Language Standards}\par\vspace{.15in}{\color{grayblue}Prepared October 2026 \quad | \quad Source edition: October 2024}\par
\end{titlepage}
\pagenumbering{roman}\tableofcontents\clearpage\pagenumbering{arabic}

\chapter*{Sources and interpretation key}
\addcontentsline{toc}{chapter}{Sources and interpretation key}
\textbf{Primary source:} ISO/IEC 9899:2024(en), \emph{Information technology --- Programming languages --- C}, fifth edition (2024). Clause and annex identifiers refer to the supplied licensed standard. This handbook contains original explanations, designs and examples; it does not reproduce the standard. The licensed standard remains the authoritative source for exact wording and conformance decisions.\par
\textbf{Secondary engineering sources:} \url{https://gcc.gnu.org/onlinedocs/gcc/Standards.html}; \url{https://clang.llvm.org/c_status.html}; \url{https://cmake.org/cmake/help/latest/prop_tgt/C_STANDARD.html}. Compiler support and diagnostics are release-specific: verify the installed compiler, target and library before relying on a feature.\par
\begin{keybox}\textbf{Interpretation categories.} \textbf{[STANDARD]} summarizes a requirement, constraint, semantic rule or defined facility and cites the applicable clause. \textbf{[OPTIONAL]} describes conditional or implementation-dependent support. \textbf{[GUIDANCE]} records an engineering recommendation, not an ISO mandate. \textbf{[TEST]} is an original check or exercise; a passing test is evidence but is not proof of full conformity.\end{keybox}
\textbf{Copyright boundary:} The supplied PDF is individually licensed. Do not redistribute it with this suite; distribute only these independently written companion materials in accordance with applicable rights.\par


\chapter{Program design}
\section{Audience and duration}
An 18-week technical curriculum for C application, systems, embedded and security engineers. The sequence assumes general programming familiarity but not mastery of C. Allow roughly six to eight hours weekly: theory, hands-on lab, assessment and peer review. The schedule is a training proposal, not an ISO prescription.
\section{Capability outcomes}
At completion, participants can read ISO/IEC 9899:2024 clauses, identify behavioral classifications, reason about lifetime and conversions, implement bounded interfaces, evaluate new C23 features, perform security reviews, and deliver a reproducible cross-toolchain build with a documented portability profile.
\chapter{Dependency graph}
\begin{keybox}\textbf{Prerequisite ordering:} Tooling/conformance $\to$ lexical/expressions $\to$ lifetimes, types and interfaces $\to$ memory/IO $\to$ modern C23 features and verification $\to$ concurrency/portability/security $\to$ design capstone $\to$ adversarial testing $\to$ final review.\end{keybox}
\section{Weekly course map}
\small\begin{longtable}{p{.09\linewidth}p{.33\linewidth}p{.21\linewidth}p{.26\linewidth}}\toprule Week & Module & Prerequisite & Submission\\\midrule\endhead
\textbf{1} & C tooling and conformance concepts & None & Environment report; first compile/log\\\hline
\textbf{2} & Lexical model and preprocessing & Week 1 & Preprocessor analysis notebook\\\hline
\textbf{3} & Scalars, promotions and expressions & Weeks 1--2 & Signedness boundary tests\\\hline
\textbf{4} & Pointers, arrays and lifetime & Week 3 & Slice implementation\\\hline
\textbf{5} & Declarations and linkage & Weeks 3--4 & Multiple-translation-unit library\\\hline
\textbf{6} & Initialization, aggregates, const & Week 5 & Data-model and init tests\\\hline
\textbf{7} & I/O and bounded parsing & Weeks 3--6 & Bounded textual parser\\\hline
\textbf{8} & Memory allocation and ownership & Weeks 4--7 & Growable buffer with failure injection\\\hline
\textbf{9} & C23 language additions & Weeks 5--8 & Feature-probed C23 micro-examples\\\hline
\textbf{10} & Integer bit operations and checked arithmetic & Weeks 3, 8--9 & Arithmetic library + boundary suite\\\hline
\textbf{11} & Binary formats, endian and encoding & Weeks 4, 7, 9 & Endian-independent record codec\\\hline
\textbf{12} & Concurrency fundamentals & Weeks 4--5, 8 & Synchronized queue\\\hline
\textbf{13} & Portability and implementation choices & Weeks 5--12 & Platform behavior register\\\hline
\textbf{14} & Security patterns and auditing & Weeks 4--13 & Code-audit report with CWE crosswalk\\\hline
\textbf{15} & Compiler verification and CI/CD & Weeks 1--14 & CMake+CI matrix and negative fixtures\\\hline
\textbf{16} & Capstone design and implementation & Weeks 1--15 & Safe record decoder and serializer\\\hline
\textbf{17} & Capstone adversarial testing & Weeks 1--16 & Fuzz/sanitizer regression bundle\\\hline
\textbf{18} & Final standard and engineering review & Weeks 1--17 & Traceability dossier + oral defense\\\hline
\bottomrule\end{longtable}\normalsize
\chapter{Weekly lesson plans and assessment}
\section{Week 1: C tooling and conformance concepts}
\textbf{Prerequisites:} None.\quad \textbf{Standard reference:} \refcl{Clauses 1--5; diagnostics, hosted/freestanding}.\par\textbf{Learning outcome:} Demonstrate standard/implementation/project policy distinctions.\par\textbf{Lab deliverable:} Environment report; first compile/log.\par\textbf{Assessment:} short concept quiz, a passing test, an intentionally failing test, a standard-versus-policy explanation, and a peer review of assumptions.

\section{Week 2: Lexical model and preprocessing}
\textbf{Prerequisites:} Week 1.\quad \textbf{Standard reference:} \refcl{5.2.1, 6.4, 6.10}.\par\textbf{Learning outcome:} Explain tokens, macro expansion and conditional compilation.\par\textbf{Lab deliverable:} Preprocessor analysis notebook.\par\textbf{Assessment:} short concept quiz, a passing test, an intentionally failing test, a standard-versus-policy explanation, and a peer review of assumptions.

\section{Week 3: Scalars, promotions and expressions}
\textbf{Prerequisites:} Weeks 1--2.\quad \textbf{Standard reference:} \refcl{6.2.5, 6.3, 6.5}.\par\textbf{Learning outcome:} Compute conversions without experimental trial and error.\par\textbf{Lab deliverable:} Signedness boundary tests.\par\textbf{Assessment:} short concept quiz, a passing test, an intentionally failing test, a standard-versus-policy explanation, and a peer review of assumptions.

\section{Week 4: Pointers, arrays and lifetime}
\textbf{Prerequisites:} Week 3.\quad \textbf{Standard reference:} \refcl{6.2.4, 6.5.7}.\par\textbf{Learning outcome:} Prove index and pointer ranges on paper and in tests.\par\textbf{Lab deliverable:} Slice implementation.\par\textbf{Assessment:} short concept quiz, a passing test, an intentionally failing test, a standard-versus-policy explanation, and a peer review of assumptions.

\section{Week 5: Declarations and linkage}
\textbf{Prerequisites:} Weeks 3--4.\quad \textbf{Standard reference:} \refcl{6.2.1--6.2.7, 6.7}.\par\textbf{Learning outcome:} Explain linkage, compatibility and header contracts.\par\textbf{Lab deliverable:} Multiple-translation-unit library.\par\textbf{Assessment:} short concept quiz, a passing test, an intentionally failing test, a standard-versus-policy explanation, and a peer review of assumptions.

\section{Week 6: Initialization, aggregates, const}
\textbf{Prerequisites:} Week 5.\quad \textbf{Standard reference:} \refcl{6.7.4, 6.7.7--6.7.12}.\par\textbf{Learning outcome:} Demonstrate complete and partial initialization semantics.\par\textbf{Lab deliverable:} Data-model and init tests.\par\textbf{Assessment:} short concept quiz, a passing test, an intentionally failing test, a standard-versus-policy explanation, and a peer review of assumptions.

\section{Week 7: I/O and bounded parsing}
\textbf{Prerequisites:} Weeks 3--6.\quad \textbf{Standard reference:} \refcl{7.23, 7.26}.\par\textbf{Learning outcome:} Reject malformed, truncated and overlong input.\par\textbf{Lab deliverable:} Bounded textual parser.\par\textbf{Assessment:} short concept quiz, a passing test, an intentionally failing test, a standard-versus-policy explanation, and a peer review of assumptions.

\section{Week 8: Memory allocation and ownership}
\textbf{Prerequisites:} Weeks 4--7.\quad \textbf{Standard reference:} \refcl{7.24.4}.\par\textbf{Learning outcome:} Retain ownership through every failure path.\par\textbf{Lab deliverable:} Growable buffer with failure injection.\par\textbf{Assessment:} short concept quiz, a passing test, an intentionally failing test, a standard-versus-policy explanation, and a peer review of assumptions.

\section{Week 9: C23 language additions}
\textbf{Prerequisites:} Weeks 5--8.\quad \textbf{Standard reference:} \refcl{6.7.10, 6.7.13, Annex M}.\par\textbf{Learning outcome:} Separate new syntax from optional target support.\par\textbf{Lab deliverable:} Feature-probed C23 micro-examples.\par\textbf{Assessment:} short concept quiz, a passing test, an intentionally failing test, a standard-versus-policy explanation, and a peer review of assumptions.

\section{Week 10: Integer bit operations and checked arithmetic}
\textbf{Prerequisites:} Weeks 3, 8--9.\quad \textbf{Standard reference:} \refcl{7.18, 7.20}.\par\textbf{Learning outcome:} Reject overflow while preserving specified output contract.\par\textbf{Lab deliverable:} Arithmetic library + boundary suite.\par\textbf{Assessment:} short concept quiz, a passing test, an intentionally failing test, a standard-versus-policy explanation, and a peer review of assumptions.

\section{Week 11: Binary formats, endian and encoding}
\textbf{Prerequisites:} Weeks 4, 7, 9.\quad \textbf{Standard reference:} \refcl{6.2.6, 6.2.9, 7.30}.\par\textbf{Learning outcome:} Round-trip bytes identically on multiple targets.\par\textbf{Lab deliverable:} Endian-independent record codec.\par\textbf{Assessment:} short concept quiz, a passing test, an intentionally failing test, a standard-versus-policy explanation, and a peer review of assumptions.

\section{Week 12: Concurrency fundamentals}
\textbf{Prerequisites:} Weeks 4--5, 8.\quad \textbf{Standard reference:} \refcl{5.1.2.4, 7.17, 7.28}.\par\textbf{Learning outcome:} Explain mutexes or release/acquire edges and test races.\par\textbf{Lab deliverable:} Synchronized queue.\par\textbf{Assessment:} short concept quiz, a passing test, an intentionally failing test, a standard-versus-policy explanation, and a peer review of assumptions.

\section{Week 13: Portability and implementation choices}
\textbf{Prerequisites:} Weeks 5--12.\quad \textbf{Standard reference:} \refcl{5.3.5, Annex J}.\par\textbf{Learning outcome:} Identify defined, unspecified and implementation-defined behavior.\par\textbf{Lab deliverable:} Platform behavior register.\par\textbf{Assessment:} short concept quiz, a passing test, an intentionally failing test, a standard-versus-policy explanation, and a peer review of assumptions.

\section{Week 14: Security patterns and auditing}
\textbf{Prerequisites:} Weeks 4--13.\quad \textbf{Standard reference:} \refcl{6.3, 6.5, 7}.\par\textbf{Learning outcome:} Prioritize actual reachable hazards and fix regressions.\par\textbf{Lab deliverable:} Code-audit report with CWE crosswalk.\par\textbf{Assessment:} short concept quiz, a passing test, an intentionally failing test, a standard-versus-policy explanation, and a peer review of assumptions.

\section{Week 15: Compiler verification and CI/CD}
\textbf{Prerequisites:} Weeks 1--14.\quad \textbf{Standard reference:} \refcl{4, 5.2.1.3, Annex K--L}.\par\textbf{Learning outcome:} Produce reproducible verification artifacts.\par\textbf{Lab deliverable:} CMake+CI matrix and negative fixtures.\par\textbf{Assessment:} short concept quiz, a passing test, an intentionally failing test, a standard-versus-policy explanation, and a peer review of assumptions.

\section{Week 16: Capstone design and implementation}
\textbf{Prerequisites:} Weeks 1--15.\quad \textbf{Standard reference:} \refcl{Clauses 4--7}.\par\textbf{Learning outcome:} Deliver all interfaces and invariants under review.\par\textbf{Lab deliverable:} Safe record decoder and serializer.\par\textbf{Assessment:} short concept quiz, a passing test, an intentionally failing test, a standard-versus-policy explanation, and a peer review of assumptions.

\section{Week 17: Capstone adversarial testing}
\textbf{Prerequisites:} Weeks 1--16.\quad \textbf{Standard reference:} \refcl{Clause families relevant to defects}.\par\textbf{Learning outcome:} Defend against malformed inputs and resource failure.\par\textbf{Lab deliverable:} Fuzz/sanitizer regression bundle.\par\textbf{Assessment:} short concept quiz, a passing test, an intentionally failing test, a standard-versus-policy explanation, and a peer review of assumptions.

\section{Week 18: Final standard and engineering review}
\textbf{Prerequisites:} Weeks 1--17.\quad \textbf{Standard reference:} \refcl{Clauses 1--7, Annexes}.\par\textbf{Learning outcome:} Explain why each claim is normative or policy and show evidence.\par\textbf{Lab deliverable:} Traceability dossier + oral defense.\par\textbf{Assessment:} short concept quiz, a passing test, an intentionally failing test, a standard-versus-policy explanation, and a peer review of assumptions.

\chapter{Required laboratory tooling}
\begin{itemize}
\item One GCC or Clang implementation with a working C23 mode; a second compiler or platform when available.
\item CMake version supporting C standard value 23, an automated test runner, code review and source control.
\item Compiler warning profile, sanitizer builds for supported platforms and static-analysis tooling.
\item Fixed example fixtures and retained failing inputs; an intentionally unsafe fixture must be marked \emph{negative-test-only}.
\end{itemize}
\chapter{Knowledge checks and practical tasks}
\section{Diagnostic classification}
Classify each scenario as language constraint, runtime undefined behavior, implementation-defined choice, unspecified behavior or project policy: incompatible redeclaration, object past lifetime, signed \texttt{char} assumptions, selected operand evaluation order, buffer naming conventions. Cite the controlling clause or record the uncertainty.
\section{Design review simulation}
Provide a proposed API \texttt{bool append(struct Buffer*, const void*, size\_t)}. Learners must specify allocation failure behavior, null treatment, aliasing, overlapping append, count overflow, borrowing/ownership, synchronization and thread safety before writing code.
\section{Migration review simulation}
Present a C17 repository with \texttt{int f()}, a zero-length \texttt{realloc}, macro named \texttt{bool}, and a platform-dependent enum serialization. Learners design a migration patch, binary compatibility test plan, and rollback strategy with clause references.
\chapter{Grading and advancement gates}
\begin{longtable}{p{.41\linewidth}r p{.41\linewidth}}\toprule Component & Weight & Pass criteria\\\midrule\endhead
Concept quizzes and clause citations & 20\% & At least 85\% accuracy and source validation.\\
Weekly lab evidence & 30\% & Tests reproducible, zero unresolved injected critical defects.\\
Portability and conformance dossier & 15\% & Traceable assumptions and a scoped claim.\\
Security assessment & 15\% & Defect proof, corrected patch and regression tests.\\
Final capstone and defense & 20\% & Interfaces, integration, adversarial tests and explanation.\\\bottomrule\end{longtable}
\guidance\ Weights and pass criteria are educational choices and do not measure ISO conformance on their own. Require remediation and re-assessment for unmet foundational safety outcomes.
\chapter{Capstone milestones}
\begin{enumerate}
\item \textbf{Week 12:} interface contract and typed wire-format specification; review adversarial input model.
\item \textbf{Week 14:} buildable decoder, serialization and resource ownership documentation.
\item \textbf{Week 16:} integration tests, compiler capability probes and conformance evidence register.
\item \textbf{Week 17:} fuzz corpus, sanitizer and static-analysis results, failure-injection closure.
\item \textbf{Week 18:} final release dossier and oral defense of normative versus engineering claims.
\end{enumerate}
\chapter{Instructor resources}
Use the Language Technical Reference for exposition, the Clause-by-Clause Guide for questions, the Migration Guide for feature adoption, the Library Reference for API preconditions, the Secure Programming Handbook for review, the Implementation Handbook for labs, and the Conformance Guide for evidence collection.
\section{Exit competency checklist}
\begin{enumerate}
\item Explains ISO requirements, optional support, undefined behavior and engineering policy without conflation.
\item Designs and reviews bounded, typed, lifetime-aware C interfaces.
\item Uses new C23 facilities only where compiler/library support is demonstrated.
\item Builds a traceable test program and records target-dependent behavior.
\item Independently justifies remediation and documents residual risks.
\end{enumerate}

\end{document}
